	\chapter{Módulos de la Plataforma}
	\label{cap:modulos}
	Este capítulo presenta, de forma breve, qué encuentra cada actor al entrar a su aplicación y cómo
	avanza una queja de un módulo al siguiente. Las funciones, reglas y estados que aquí se mencionan
	están especificados en el Capítulo~\ref{cap:analisis_diseno}; entre paréntesis se indica el caso de
	uso que implementa cada pantalla.

	\section{Portal del quejoso}
	\label{sec:modulo_quejoso}
	Es la única aplicación dirigida al público y se publica en la raíz del dominio. Su área pública
	permite presentar una queja con un formulario en pasos numerados, que previsualiza los archivos y
	señala en cada campo la regla que incumple (CU-QJ-07); consultar una queja con el folio y el correo
	(CU-QJ-11); activar la cuenta de seguimiento, iniciar sesión y recuperar la contraseña (CU-QJ-01 a
	CU-QJ-04); responder sin sesión las citas y los acuerdos recibidos por correo (CU-PC-04, CU-QJ-14), y
	consultar el asistente de preguntas frecuentes (CU-QJ-06). El panel autenticado reúne el resumen de
	las quejas propias, su detalle con la línea de tiempo y las acciones que permite cada estatus
	---editar o cancelar en \texttt{RECIBIDA}, corregir en \texttt{RECHAZADA}--- (CU-QJ-08 a CU-QJ-10,
	CU-QJ-12, CU-QJ-15), los acuerdos de conciliación (CU-QJ-13, CU-QJ-14), las notificaciones
	(CU-NOT-01) y el perfil (CU-QJ-05).

	\section{Respuesta de la persona señalada}
	\label{sec:modulo_denunciado}
	La página \texttt{/denunciado/respuesta} forma parte del portal, pero está dirigida a la persona a
	quien se atribuyen los hechos. Muestra el folio de la queja recibido en el enlace, solicita su
	identidad, su versión de los hechos, su credencial y sus evidencias, exige su aviso de privacidad y
	entrega un folio de respuesta (RF-D01, RF-D02).

	\section{Panel de Recepción}
	\label{sec:modulo_revision}
	Su pantalla de inicio de sesión es la entrada única del personal, que es dirigido al panel de su rol
	(CU-ADM-07). La bandeja muestra los indicadores del día, la cola paginada y la actividad de cada
	recepcionista (CU-REC-01). Al abrir una queja, esta queda asignada a quien la abrió y se presenta la
	pantalla de validación: los datos y documentos a la izquierda (CU-REC-03) y, a la derecha, la lista
	de verificación de los cinco requisitos de admisión con su fundamento, el catálogo de documentos
	aceptados y el cómputo automático del plazo de 90 días (CU-REC-02). Con los cinco requisitos
	satisfechos se habilita el turnado, que muestra los antecedentes de la queja y sugiere al abogado de
	menor carga (CU-REC-05); con alguno incumplido, solo puede regresarse al quejoso con observaciones
	(CU-REC-04). El panel incluye además el registro en ventanilla, que reutiliza el formulario del
	portal (CU-REC-08), y el historial con exportación a Excel (CU-REC-07).

	\section{Panel de Primer Contacto}
	\label{sec:modulo_primer_contacto}
	La bandeja lista los expedientes turnados con su prioridad y estatus (CU-PC-02). Al abrir un
	expediente, este pasa a \texttt{EN\_ANALISIS} (CU-PC-03) y se presenta en cuatro pestañas:

	\begin{itemize}
		\item \textbf{Análisis IA.} Borrador de resumen con botones para copiarlo y ver la narración
		completa; fragmentos clave etiquetados por tema ---inicio del contexto, hechos principales,
		consecuencia, trato recibido---; los cuatro antecedentes más parecidos con su porcentaje de
		similitud y una gráfica de barras; y el acceso al análisis detallado de antecedentes, donde pueden
		guardarse o buscarse manualmente. Si el componente no responde, la pestaña lo indica.
		\item \textbf{Expediente.} Narración, datos de las partes, evidencias e historial del expediente
		como línea de tiempo.
		\item \textbf{Gestión.} Oficios con su PDF (CU-PC-05), solicitud y versión de la persona señalada
		(RF-D03), citas (CU-PC-04) y propuesta de conciliación (CU-PC-10).
		\item \textbf{Resolución.} Dictamen de competencia, con el responsable en Subdefensoría elegido
		por carga, o de improcedencia con su remisión (CU-PC-06 a CU-PC-09).
	\end{itemize}

	La agenda presenta las citas en calendario y en lista, y su ventana de programación muestra cada
	horario como disponible, ocupado o no disponible (CU-PC-04).

	\section{Panel de Subdefensoría}
	\label{sec:modulo_subdefensoria}
	El panel se organiza alrededor del siguiente paso de cada expediente. El inicio muestra los
	pendientes del día ordenados por urgencia y los expedientes que vencen pronto; ``Mis expedientes''
	los agrupa por etapa con una barra de plazo (CU-SD-02, CU-SD-03); el expediente reúne lo actuado en
	Primer Contacto, en modo de solo lectura, y las acciones de investigación ---oficios, recordatorios,
	respuesta de la unidad y acuerdos--- (CU-SD-05 a CU-SD-09); y ``Plazos y oficios'' concentra las
	alertas de vencimiento (CU-SD-04).

	\section{Panel de la Defensora}
	\label{sec:modulo_admin}
	Es exclusivo del rol \texttt{DEFENSOR}. En la sección de supervisión, el panel general muestra los
	indicadores del periodo, la distribución de las quejas por etapa, su evolución mensual, las unidades
	con más quejas, la carga del personal y los casos que requieren atención según las reglas de la
	Tabla~\ref{tab:reglas_atencion}; cada número abre la lista filtrada que lo explica. Desde el
	explorador de quejas, la Defensora filtra por prioridad y revisa directamente el detalle de las
	quejas prioritarias (CU-ADM-03). En la sección de administración gestiona el personal (CU-ADM-01,
	CU-ADM-02), el catálogo de dependencias (CU-ADM-06), las plantillas oficiales (CU-ADM-04), las
	preguntas frecuentes (CU-BOT-01) y los respaldos con la bitácora (CU-ADM-05).

	\rowcolors{2}{rowblue}{white}
	\small
	\renewcommand{\arraystretch}{1.2}
	\setlength{\tabcolsep}{4pt}
	\begin{longtable}{|p{5.0cm}|p{9.5cm}|}
		\caption{Reglas de ``Casos que requieren atención''}
		\label{tab:reglas_atencion} \\
		\hline
		\rowcolor{headerblue}
		\headercell{Regla} & \headercell{Condición} \\
		\hline
		\endfirsthead
		\hline
		\rowcolor{headerblue}
		\headercell{Regla} & \headercell{Condición} \\
		\hline
		\endhead
		Sin tomar & Más de 3 días en Recepción sin abrirse \\ \hline
		Corregida sin revisar & Más de 3 días desde la corrección \\ \hline
		Validación detenida & Abierta por una recepcionista hace más de 1 día \\ \hline
		Observaciones pendientes & El quejoso no ha corregido en más de 5 días \\ \hline
		Sin abogado & En Primer Contacto sin abogado asignado \\ \hline
		Vencida (Primer Contacto) & Más de 10 días en Primer Contacto \\ \hline
		Sin turnar a Subdefensoría & Procedente hace más de 15 días \\ \hline
		Vencida (seguimiento) & Más de 60 días en Subdefensoría \\ \hline
	\end{longtable}
	\normalsize

	Los días se muestran en rojo cuando alcanzan el doble del plazo de la regla.
